\section{Теоретический анализ методов обнаружения аномалий и обоснование требований к системе}
\label{ch:theory}

\noindent
\textit{Содержание главы соответствует первому этапу постановки
задачи (см.~ВКР, Этап~1): изучение литературы и существующих методов
обнаружения аномалий в сетевом трафике; анализ систем IDS/IPS и их
компонентов; обзор методов создания ИИ-злоумышленника; формулировка
гипотезы исследования и требований к системе.}\\

\subsection{Корпоративная сеть как объект мониторинга}

Под корпоративной сетью в настоящей работе понимается совокупность
взаимосвязанных программно-аппаратных средств, обеспечивающих обмен
данными между рабочими станциями, серверами, сетевым оборудованием и
внешними сегментами. С точки зрения мониторинга безопасности
корпоративная сеть характеризуется тремя свойствами, существенными
для построения NIDS:

\begin{enumerate}
    \item высокой неоднородностью трафика, приводящей к мультимодальному
    распределению признаков и затрудняющей применение единых пороговых
    правил~\cite{khraisat2019survey, ahmad2021nids};
    \item доминированием легитимного трафика, что создаёт сильный
    дисбаланс классов и требует специализированных методов обучения
    и оценки~\cite{sommer2010outside, chawla2002smote, lin2017focal};
    \item динамичностью характеристик, приводящей к явлению дрейфа
    распределений~\cite{gama2014drift, andresini2021insomnia}.
\end{enumerate}

Объектами мониторинга в работе принимаются: сетевые соединения
(flow) на уровне L3/L4, сессии прикладного уровня и временные
последовательности соединений в скользящих окнах. Это согласуется
с подходом, принятым в большинстве современных систем сетевого
анализа трафика и в открытых
датасетах~\cite{moustafa2015unsw, sharafaldin2018toward}.

В качестве рабочей таксономии угроз использована классификация
датасета UNSW-NB15, выделяющая девять семейств атакующего трафика:
\emph{Fuzzers, Analysis, Backdoors, DoS, Exploits, Generic,
Reconnaissance, Shellcode, Worms}~\cite{moustafa2015unsw}.

\subsection{Классификация и принципы построения IDS}

В соответствии с устоявшейся в литературе типологией IDS делятся на
три класса~\cite{khraisat2019survey, ahmad2021nids, ferrag2020dlcyber}:
\textbf{сигнатурные} (совпадение с заранее заданным шаблоном),
\textbf{аномальные} (отклонение от модели нормы) и \textbf{гибридные}.
Сигнатурный класс обеспечивает высокую точность по известным угрозам,
но не способен обнаруживать ранее не встречавшиеся атаки (\eng{zero-day}).
Аномальный класс расширяет покрытие на неизвестные сценарии за счёт
повышенной частоты ложных срабатываний и чувствительности к дрейфу.

С точки зрения исходного представления данных различают \emph{packet-based}
(анализ полезной нагрузки, DPI) и \emph{flow-based} анализ.
Стандартами экспорта flow-записей являются Cisco NetFlow~\cite{rfc3954} и
IETF IPFIX~\cite{rfc7011}. В настоящей работе принят flow-based подход
по причине его совместимости с шифрованным трафиком, ресурсной
эффективности и прямой переносимости на инфраструктуры реальных
корпоративных сетей.

\subsection{Методы обнаружения аномалий}

Постановка задачи: для пространства признаков \(X \subseteq \mathbb{R}^d\)
требуется построить функцию скоринга \(s\colon X \to \mathbb{R}\) и
порог \(\tau\) такие, что
\begin{equation}
\hat{y}(x) = \mathbb{I}\bigl[s(x) > \tau\bigr],
\quad x \in X,
\label{eq:detector}
\end{equation}
где \(\hat{y}(x) = 1\) интерпретируется как ``аномалия / атака''.

\textbf{Статистические и энтропийные методы.} Простейший подход ---
\(z\)-оценка и многомерное расстояние Махаланобиса:
\begin{equation}
d_M(x) = \sqrt{(x-\mu)^{\top}\Sigma^{-1}(x-\mu)}.
\label{eq:mahalanobis}
\end{equation}
Энтропия Шеннона \(H(P) = -\sum_a P(a)\log P(a)\) используется для
оценки распределений IP-адресов и портов в скользящем окне.

\textbf{Isolation Forest}~\cite{liu2008isolation} определяет
скоринг через ожидаемую глубину изоляции точки:
\begin{equation}
s(x, n) = 2^{-\frac{\mathbb{E}[h(x)]}{c(n)}},\quad
c(n) \approx 2(\ln(n-1) + 0.5772) - \frac{2(n-1)}{n}.
\label{eq:iforest}
\end{equation}
Метод имеет линейную сложность и хорошо подходит для unsupervised-режима.

\textbf{One-Class SVM}~\cite{scholkopf2001ocsvm} ищет в спрямляющем
пространстве минимальную гиперсферу, содержащую долю \(1-\nu\)
обучающих наблюдений; параметр \(\nu \in (0,1]\) задаёт верхнюю
оценку доли выбросов.

\subsection{Глубинное обучение для анализа flow-данных}

В работе рассмотрено восемь типов нейросетевых архитектур.

\textbf{MLP} реализуется рекурсией
\(h^{(l)} = \phi(W^{(l)} h^{(l-1)} + b^{(l)})\) и служит нижней
границей DL-моделей.

\textbf{1D-CNN} применяет операцию свёртки
\((h * w)_t = \sum_{k=0}^{K-1} w_k h_{t-k}\) для извлечения локальных
паттернов в скользящем окне.

\textbf{LSTM}~\cite{hochreiter1997lstm} вводит ячейку состояния
\(c_t\) и три гейта:
\begin{align}
i_t &= \sigma(W_i [h_{t-1}, x_t] + b_i), \nonumber\\
f_t &= \sigma(W_f [h_{t-1}, x_t] + b_f), \nonumber\\
o_t &= \sigma(W_o [h_{t-1}, x_t] + b_o), \nonumber\\
\tilde{c}_t &= \tanh(W_c [h_{t-1}, x_t] + b_c), \nonumber\\
c_t &= f_t \odot c_{t-1} + i_t \odot \tilde{c}_t, \nonumber\\
h_t &= o_t \odot \tanh(c_t).
\label{eq:lstm}
\end{align}
Гейтовая структура обеспечивает устойчивое распространение градиента и
способность хранить долгосрочный контекст --- свойство, востребованное
при обнаружении медленных атак.

\textbf{GRU}~\cite{cho2014gru} --- упрощённая RNN с двумя гейтами;
\textbf{BiLSTM} использует прямой и обратный LSTM для двунаправленного
контекста.

\textbf{Гибрид CNN-LSTM} применяет последовательное преобразование
\[
x_t \xrightarrow{\text{Conv1D}} z_t \xrightarrow{\text{LSTM}} h_t
\xrightarrow{\text{FC}} \hat{p}_t,
\]
сочетая локальные паттерны и временные зависимости.

\textbf{TCN}~\cite{bai2018tcn} использует причинные дилатированные
свёртки с экспоненциально растущим рецептивным полем.

\textbf{Transformer-encoder}~\cite{vaswani2017attention} основан на
само-внимании:
\begin{equation}
\operatorname{Attention}(Q, K, V) =
\operatorname{softmax}\!\left(\frac{QK^{\top}}{\sqrt{d_k}}\right) V.
\label{eq:attention}
\end{equation}

\textbf{Autoencoder и VAE.} Автоэнкодер минимизирует ошибку
реконструкции \(\|x - D(E(x))\|^2\); вариационный
автоэнкодер~\cite{kingma2014vae} обучает вероятностное кодирование
максимизацией ELBO:
\begin{equation}
\mathcal{L}_{\text{VAE}} =
\mathbb{E}_{q(z|x)}[\log p(x|z)] - \KL(q(z|x)\|p(z)).
\label{eq:vae}
\end{equation}

Качественное сравнение архитектур по характеристикам приведено в
\tabref{tab:dl-arch-compare}.

\begin{table}[H]
\centering
\caption{Качественное сравнение нейросетевых архитектур для NIDS}
\label{tab:dl-arch-compare}
\renewcommand{\arraystretch}{1.2}
\begin{tabularx}{\linewidth}{@{}lXXXX@{}}
\toprule
\textbf{Модель} & \textbf{Локальные паттерны} & \textbf{Временные зависимости} &
\textbf{Сложность обучения} & \textbf{Inference latency} \\
\midrule
MLP        & нет          & нет        & низкая  & низкая \\
1D-CNN     & да           & ограничено & средняя & низкая \\
LSTM       & нет          & да         & средняя & средняя \\
GRU        & нет          & да         & средняя & низкая \\
BiLSTM     & нет          & двунаправ. & высокая & средняя \\
CNN-LSTM   & да           & да         & высокая & средняя \\
TCN        & да           & фиксиров.  & средняя & низкая \\
Transformer & через attention & да     & высокая & средняя \\
AE / VAE   & косвенно     & нет        & средняя & низкая \\
\bottomrule
\end{tabularx}
\end{table}

\subsection{Моделирование атак с использованием ИИ}

В литературе выделяют три семейства подходов к формированию атакующего
трафика для тестирования NIDS:

\begin{enumerate}
    \item \textbf{Шаблонная (\eng{rule-based, template-based}) генерация}:
    атаки задаются параметрическими шаблонами потоков;
    применение --- от классических работ DARPA/KDD до утилит
    \eng{traffic replay}~\cite{lippmann2000darpa, gharib2016evalframework}.
    \item \textbf{Генеративные модели}: GAN, cGAN, WGAN, VAE для синтеза
    реалистичных flow-векторов и последовательностей. Эталонная
    работа --- Ring и соавт.~\cite{ring2019flowgan}; вариационные подходы ---
    \cite{lopez2017cvae}.
    \item \textbf{Adversarial-агенты на обучении с подкреплением}:
    агент обучается формировать последовательности действий,
    обходящих детектор~\cite{apruzzese2020adversarial}.
\end{enumerate}

Классическая постановка GAN~\cite{goodfellow2014gan}:
\begin{equation}
\min_{G}\max_{D}
V(D,G) = \mathbb{E}_{x \sim p_{\text{data}}}[\log D(x)] +
         \mathbb{E}_{z \sim p_z}[\log(1 - D(G(z)))].
\label{eq:gan}
\end{equation}
Conditional GAN~\cite{mirza2014cgan} вводит обусловливающую переменную
\(c\) (тип атаки); WGAN-GP~\cite{arjovsky2017wgan, gulrajani2017wgangp}
заменяет JS-расстояние Wasserstein-1 и обеспечивает 1-липшицевость
дискриминатора через градиентный штраф.

В рамках настоящей работы обоснован выбор \textbf{гибридной шаблонной
архитектуры} ИИ-злоумышленника:
\[
\text{Templates} + \text{FSM-планировщик} + \text{Drift-инжектор}
\]
без cGAN-составляющей. Аргументы: высокая воспроизводимость и
прозрачная ground-truth, семантическая согласованность шаблонов
с таксономией UNSW-NB15, отсутствие риска коллапса режима,
интерпретируемость для аналитика SOC. cGAN сохраняется в обзоре
как теоретическая альтернатива и направление дальнейшего развития.

\subsection{Сравнительный анализ датасетов NIDS}

В \tabref{tab:datasets-compare} приведена сводная экспертная оценка
десяти публично доступных датасетов NIDS по восьми
критериям, систематизированным в~\cite{gharib2016evalframework,
sharafaldin2018toward}.

\begin{table}[H]
\centering
\caption{Сводное сравнение публично доступных датасетов NIDS}
\label{tab:datasets-compare}
\renewcommand{\arraystretch}{1.2}
\setlength{\tabcolsep}{3pt}
\small
\begin{tabularx}{\linewidth}{@{}lcccccccX@{}}
\toprule
\textbf{Датасет} & \textbf{K1} & \textbf{K2} & \textbf{K3} & \textbf{K4} &
\textbf{K5} & \textbf{K6} & \textbf{K7} & \textbf{K8: ограничения} \\
\midrule
KDDCup99~\cite{tavallaee2009nslkdd}            & $-$    & $-$ & $-$ & $\circ$ & $\circ$ & $-$    & $+$ & устаревшая модель угроз \\
NSL-KDD~\cite{tavallaee2009nslkdd}             & $-$    & $-$ & $\circ$ & $+$ & $\circ$ & $-$    & $+$ & таксономия 1998~г.; малый объём для DL \\
\textbf{UNSW-NB15}~\cite{moustafa2015unsw}     & $+$    & $+$ & $+$ & $\circ$ & $+$    & $+$    & $+$ & доминирование Generic \\
CICIDS2017~\cite{sharafaldin2018toward}        & $+$    & $+$ & $\circ$ & $\circ$ & $+$    & $+$    & $+$ & дефекты разметки; дубликаты \\
CSE-CIC-IDS2018~\cite{sharafaldin2018toward}   & $+$    & $+$ & $\circ$ & $\circ$ & $+$    & $+$    & $\circ$ & унаследованные дефекты \\
TON\_IoT~\cite{moustafa2021toniot}             & $+$    & $\circ$ & $+$ & $\circ$ & $+$    & $+$    & $+$ & ориентация на IoT \\
Bot-IoT~\cite{koroniotis2019botiot}            & $\circ$& $-$ & $+$ & $-$    & $+$    & $+$    & $+$ & доля атак $>97\%$ \\
Kyoto 2006+~\cite{song2011kyoto}               & $\circ$& $\circ$ & $\circ$ & $\circ$ & $\circ$ & $\circ$ & $\circ$ & honeypot-смещение \\
DARPA~\cite{lippmann2000darpa}                 & $-$    & $-$ & $\circ$ & $\circ$ & $-$    & $\circ$ & $+$ & 1998~г.; искусственный трафик \\
LITNET-2020~\cite{damasevicius2020litnet}      & $+$    & $+$ & $\circ$ & $\circ$ & $+$    & $+$    & $\circ$ & таксономический дисбаланс \\
\bottomrule
\end{tabularx}
\end{table}

\noindent
K1 --- актуальность; K2 --- репрезентативность нормы;
K3 --- качество разметки; K4 --- дисбаланс классов;
K5 --- пригодность для DL; K6 --- flow-формат;
K7 --- воспроизводимость; K8 --- ограничения.
``$+$'' --- соответствует, ``$\circ$'' --- частично, ``$-$'' --- нет.

В работе принят датасет \textbf{UNSW-NB15} как основной (наибольшее
покрытие критериев), \textbf{CICIDS2017} --- как вспомогательный для
кросс-валидации.

\subsection{Метрики качества}

Для бинарной задачи введены матрица ошибок (\(TP, TN, FP, FN\)) и
производные метрики:
\begin{align}
\text{Precision} &= \frac{TP}{TP + FP},
\quad
\text{Recall} = \frac{TP}{TP + FN},
\label{eq:pr}\\
\text{F1} &= \frac{2\,\text{Precision}\,\text{Recall}}
                 {\text{Precision} + \text{Recall}},
\quad
\text{FPR} = \frac{FP}{FP + TN}.\label{eq:f1}
\end{align}

\textbf{MCC} (Matthews correlation coefficient) рекомендуется при
сильном дисбалансе:
\begin{equation}
\text{MCC} =
\frac{TP \cdot TN - FP \cdot FN}
{\sqrt{(TP+FP)(TP+FN)(TN+FP)(TN+FN)}}.
\label{eq:mcc}
\end{equation}

\textbf{ROC-AUC} и \textbf{PR-AUC} интегрируют качество модели по всем
порогам; PR-AUC более чувствителен к редкому классу и принят как
основная интегральная метрика. К эксплуатационным метрикам относятся
\emph{time-to-detect} (TTD), \emph{inference latency}, \emph{throughput}
и \emph{FP-per-time}.

\textbf{Температурное масштабирование}~\cite{guo2017calibration}
согласует выходную вероятность с эмпирической частотой:
\begin{equation}
\hat{p}_{T}(y=1\mid x) = \sigma\!\left(\frac{z(x)}{T}\right),
\label{eq:temp-scaling}
\end{equation}
где \(z(x)\) --- логит, \(T\) --- параметр, подбираемый на
валидационной выборке. Качество калибровки оценивается по ECE
(\eng{Expected Calibration Error}).

\textbf{Метрики дрейфа.} PSI (Population Stability Index),
KL-расхождение и MMD (Maximum Mean Discrepancy):
\begin{align}
\PSI &= \sum_{b}(p_b - q_b)\ln\frac{p_b}{q_b},
\label{eq:psi}\\
\MMD^{2}(P, Q) &= \mathbb{E}_{P,P}[k(x,x')] - 2\mathbb{E}_{P,Q}[k(x,y)]
                + \mathbb{E}_{Q,Q}[k(y,y')].\label{eq:mmd}
\end{align}

\subsection{Формализация требований и гипотеза}

Концептуальная модель системы состоит из трёх подсистем:
предобработки, детекции и активного тестирования
(ИИ-злоумышленник). Сформированы 10 функциональных и 9 нефункциональных
требований; ключевые количественные пороги приведены в
\tabref{tab:requirements}.

\begin{table}[H]
\centering
\caption{Ключевые нефункциональные требования к системе}
\label{tab:requirements}
\renewcommand{\arraystretch}{1.2}
\begin{tabularx}{\linewidth}{@{}lXl@{}}
\toprule
\textbf{ID} & \textbf{Содержание} & \textbf{Порог} \\
\midrule
NFR-1 & F1 на UNSW-NB15 (binary) & $\geq 0{,}90$ \\
NFR-1 & PR-AUC на UNSW-NB15      & $\geq 0{,}95$ \\
NFR-2 & FPR при Recall $\geq 0{,}85$ & $\leq 0{,}01$ \\
NFR-3 & Inference latency (CPU, batch=1) & $\leq 50$ мс/окно \\
NFR-4 & Throughput (CPU, батч)           & $\geq 1000$ EPS \\
NFR-5 & Time-to-detect (DDoS / Scan / Brute) & $\leq 1$ окна \\
NFR-6 & $\Delta\text{F1}$ при PSI $\leq 0{,}25$ & $\leq 0{,}10$ \\
NFR-7 & Перенос UNSW $\to$ CICIDS2017 без дообуч. ($\Delta\text{F1}$) & $\leq 0{,}20$ \\
NFR-8 & Воспроизводимость одной командой & да \\
NFR-9 & Стабильность $\sigma_{\text{F1}}$ (\(n \geq 5\) seed) & $\leq 0{,}01$ \\
\bottomrule
\end{tabularx}
\end{table}

Формально сформулированы три гипотезы исследования:
основная \(H_0^{\text{main}}\) о превосходстве CNN-LSTM над ансамблем
baselines на UNSW-NB15; дополнительная \(H_0^{\text{aux}}\) о роли
ИИ-злоумышленника в воспроизводимом тестировании; гипотеза
переносимости \(H_0^{\text{transfer}}\) на CICIDS2017. Проверка
проводится на этапе~\ref{ch:experiments}.

\subsection*{Выводы по главе 1}

\begin{enumerate}
    \item Корпоративная сеть как объект мониторинга характеризуется
    неоднородностью трафика, дисбалансом классов и динамичностью
    распределений; обоснован выбор flow-уровня анализа,
    совместимого с NetFlow/IPFIX.
    \item Сигнатурные, аномальные и гибридные подходы NIDS
    дополняют, а не заменяют друг друга; в работе принят аномальный
    нейросетевой подход.
    \item Рассмотрены восемь типов нейросетевых архитектур, для каждой
    приведены математические основы; гибрид CNN-LSTM покрывает оба
    измерения структуры flow-данных и принимается в качестве основной
    proposed-архитектуры.
    \item Систематизированы три семейства подходов к моделированию
    атак; обоснован выбор гибридной шаблонной архитектуры
    ИИ-злоумышленника без cGAN-составляющей.
    \item Из десяти датасетов выбран UNSW-NB15 как основной и
    CICIDS2017 как вспомогательный.
    \item Сформированы количественные требования и гипотезы исследования.
\end{enumerate}
